\documentclass[10pt,a4paper]{amsart}
\usepackage{amsmath,amssymb,amsthm,graphicx}
% Retained preamble for the Stacks Project chapter morphisms.tex (Morphisms of Schemes).
% The chapter file pinned-original-source.tex carries only \input{preamble} on its first line, so no
% preamble text lives inside the chapter itself.  The Stacks Project uses one single shared preamble for
% every chapter; the snapshot stages it verbatim in sources/stacks-chow/preamble.tex.  The block below is
% copied mechanically, line for line, from that file's lines 14-18 (the Xy-pic setup) and lines 154-217
% (the native theorem declarations together with the mathematical macros).  Nothing was added, removed or
% redefined.  The wrapper class and the amsmath/amssymb/amsthm/graphicx base packages are supplied by the
% assembler, so the class/\externaldocument part of the shared preamble is not repeated here.
% For commutative diagrams we use Xy-pic
\usepackage[all]{xy}

% We use 2cell for 2-commutative diagrams.
\xyoption{2cell}
% Theorem environments.
%
\theoremstyle{plain}
\newtheorem{theorem}[subsection]{Theorem}
\newtheorem{proposition}[subsection]{Proposition}
\newtheorem{lemma}[subsection]{Lemma}

\theoremstyle{definition}
\newtheorem{definition}[subsection]{Definition}
\newtheorem{example}[subsection]{Example}
\newtheorem{exercise}[subsection]{Exercise}
\newtheorem{situation}[subsection]{Situation}

\theoremstyle{remark}
\newtheorem{remark}[subsection]{Remark}
\newtheorem{remarks}[subsection]{Remarks}

\numberwithin{equation}{subsection}

% Macros
%
\def\lim{\mathop{\mathrm{lim}}\nolimits}
\def\colim{\mathop{\mathrm{colim}}\nolimits}
\def\Spec{\mathop{\mathrm{Spec}}}
\def\Hom{\mathop{\mathrm{Hom}}\nolimits}
\def\Ext{\mathop{\mathrm{Ext}}\nolimits}
\def\SheafHom{\mathop{\mathcal{H}\!\mathit{om}}\nolimits}
\def\SheafExt{\mathop{\mathcal{E}\!\mathit{xt}}\nolimits}
\def\Sch{\mathit{Sch}}
\def\Mor{\mathop{\mathrm{Mor}}\nolimits}
\def\Ob{\mathop{\mathrm{Ob}}\nolimits}
\def\Sh{\mathop{\mathit{Sh}}\nolimits}
\def\NL{\mathop{N\!L}\nolimits}
\def\CH{\mathop{\mathrm{CH}}\nolimits}
\def\proetale{{pro\text{-}\acute{e}tale}}
\def\etale{{\acute{e}tale}}
\def\QCoh{\mathit{QCoh}}
\def\Ker{\mathop{\mathrm{Ker}}}
\def\Im{\mathop{\mathrm{Im}}}
\def\Coker{\mathop{\mathrm{Coker}}}
\def\Coim{\mathop{\mathrm{Coim}}}

% Boxtimes
%
\DeclareMathSymbol{\boxtimes}{\mathbin}{AMSa}{"02}

%
% Macros for moduli stacks/spaces
%
\def\QCohstack{\mathcal{QC}\!\mathit{oh}}
\def\Cohstack{\mathcal{C}\!\mathit{oh}}
\def\Spacesstack{\mathcal{S}\!\mathit{paces}}
\def\Quotfunctor{\mathrm{Quot}}
\def\Hilbfunctor{\mathrm{Hilb}}
\def\Curvesstack{\mathcal{C}\!\mathit{urves}}
\def\Polarizedstack{\mathcal{P}\!\mathit{olarized}}
\def\Complexesstack{\mathcal{C}\!\mathit{omplexes}}
% \Pic is the operator that assigns to X its picard group, usage \Pic(X)
% \Picardstack_{X/B} denotes the Picard stack of X over B
% \Picardfunctor_{X/B} denotes the Picard functor of X over B
\def\Pic{\mathop{\mathrm{Pic}}\nolimits}
\def\Picardstack{\mathcal{P}\!\mathit{ic}}
\def\Picardfunctor{\mathrm{Pic}}
\def\Deformationcategory{\mathcal{D}\!\mathit{ef}}

% --- appended 2026-10-12 (Chief): the project own definitions of the reference/slogan/history
% display environments, copied mechanically from the shared Stacks preamble lines 8-12
% (\usepackage{verbatim} plus the three \newenvironment lines, which the chapter preamble keeps
% verbatim from the project preamble). Nothing added, removed or redefined: these environments
% carry no mathematical content and the project itself typesets them as comments. Without these
% three lines a counted statement whose own text uses them fails with Environment slogan undefined.
\usepackage{verbatim}
\newenvironment{reference}{\comment}{\endcomment}
\newenvironment{slogan}{\comment}{\endcomment}
\newenvironment{history}{\comment}{\endcomment}
\usepackage{cleveref}
\crefname{originallemma}{original lemma}{original lemmas}
\Crefname{originallemma}{Original lemma}{Original lemmas}
\makeatletter
\newlabel{schemes-lemma-characterize-points}{{01J9}{source}{}{sourcepointer-1}{}}
\newlabel{schemes-lemma-characterize-points@cref}{{[originallemma][0][]01J9}{[1][0][]source}{}{}{}}
\makeatother
\crefname{originallemma}{original lemma}{original lemmas}
\Crefname{originallemma}{Original lemma}{Original lemmas}
\makeatletter
\newlabel{schemes-lemma-points-fibre-product}{{01JT}{source}{}{sourcepointer-2}{}}
\newlabel{schemes-lemma-points-fibre-product@cref}{{[originallemma][0][]01JT}{[1][0][]source}{}{}{}}
\makeatother
\title{Four equivalent characterizations of universally injective morphisms, including surjectivity of the diagonal}
\author{The Stacks Project Authors}
\date{stacks-project:morphisms.tex:local-snapshot-20261012; GNU FDL 1.2 or later}
\begin{document}
\maketitle
\noindent\textit{Editorial note.} All delivered mathematical text is the author's own, copied from the pinned Stacks Project chapter file pinned-original-source.tex; no statement or proof step was rewritten, reordered or removed.\par\medskip
\noindent\textit{Editorial note.} Following the multi-claim convention for this source, source\_url is the source-level URL https://stacks.math.columbia.edu/ and the source label lemma-universally-injective is carried in claim\_key and here rather than in a per-claim tag URL.\par\medskip
\noindent\textit{Editorial note.} The preamble is the chapter's shared Stacks Project preamble, retained as sources/stacks-morphisms/retained-preamble.tex: a line-for-line excerpt of the shared Stacks preamble staged with this snapshot. No macro was invented or redefined, and this card adds no preamble text of its own.\par\medskip
\noindent\textit{Editorial note.} Statement span 1577-1593 and proof span 1595-1667 come from the halfA pool card operations/pool-cards/stacks-morphisms-halfA.json and were re-read in the pinned chapter file before this card was written.\par\medskip
\noindent\textit{Editorial note.} The counted statement contains a reference environment holding the author's citation of EGA; the project's own shared preamble (retained in full at 88 lines after the 2026-10-12 repair) defines that environment as a comment environment, so no bibliography is attached and none is needed.\par\medskip
\noindent\textit{Editorial note.} The statement also contains an xymatrix diagram of fields inside the proof; it is copied verbatim and the retained preamble loads Xy-pic.\par\medskip
\noindent\textit{Editorial note.} 2 cross-references occur inside the counted spans; they are listed in external\_references as navigation pointers only.\par\medskip
\noindent\textit{Allowed standard prerequisites (editorial).}\begin{itemize}
\item Points of a scheme over a field and of a fibre product: morphisms from the spectrum of a field to X correspond to pairs consisting of a point and an embedding of its residue field, and points of X times\_S X correspond to quadruples with a prime ideal in the tensor product of the two residue fields. Cited as a schemes chapter lemma at proof lines 1826, 1848, 1858 and 1872-1874. It is the dictionary that converts the morphism-counting condition (1) and the diagonal condition (4) into the pointwise statements (2) and (3).
\item Fibre products and base change of the diagonal, and the definition of universally injective as base change being injective on points. Used at proof lines 1813-1818 and 1879-1883. Identifies sections of a base change with points of the fibre and turns injectivity of base changes into the diagonal being a section of an injective morphism.
\item Purely inseparable field extensions in terms of uniqueness of embeddings into a common extension field. Used at proof lines 1841-1846 and 1877-1879. Supplies the field-theoretic step that two different embeddings produce two different morphisms from the spectrum of a field, and conversely.
\end{itemize}\medskip
\ifdefined\hypertarget\hypertarget{sourcepointer-1}{}\fi
\noindent\textit{Source reference.} Original lemma 01J9: focused-seeds/AGP/sources/stacks-snapshot-20261011/tags.txt, source lines 2006--2006. Same-chapter reference used in the retained author text at source line(s) 1612, 1629, 1638. It is supplied only as a navigation pointer so that the retained text compiles; the counted item delivers the author's own statement and the author's complete proof, and no part of the referenced item is copied in or substituted for the counted argument.\par\medskip
\ifdefined\hypertarget\hypertarget{sourcepointer-2}{}\fi
\noindent\textit{Source reference.} Original lemma 01JT: focused-seeds/AGP/sources/stacks-snapshot-20261011/tags.txt, source lines 2026--2026. Same-chapter reference used in the retained author text at source line(s) 1326, 1485, 1652, 3643, 10040. It is supplied only as a navigation pointer so that the retained text compiles; the counted item delivers the author's own statement and the author's complete proof, and no part of the referenced item is copied in or substituted for the counted argument.\par\medskip
\begin{lemma}
\label{lemma-universally-injective}
\begin{reference}
\cite[Proposition 3.7.1]{EGA1-second}
\end{reference}
Let $f : X \to S$ be a morphism of schemes.
The following are equivalent:
\begin{enumerate}
\item For every field $K$ the induced map
$\Mor(\Spec(K), X) \to \Mor(\Spec(K), S)$
is injective.
\item The morphism $f$ is universally injective.
\item The morphism $f$ is radicial.
\item The diagonal morphism $\Delta_{X/S} : X \longrightarrow X \times_S X$
is surjective.
\end{enumerate}
\end{lemma}

\begin{proof}
Let $K$ be a field, and let $s : \Spec(K) \to S$ be a morphism.
Giving a morphism $x : \Spec(K) \to X$ such that $f \circ x = s$
is the same as giving a section of the projection
$X_K = \Spec(K) \times_S X \to \Spec(K)$, which in turn
is the same as giving a point $x \in X_K$ whose residue field is $K$.
Hence we see that (2) implies (1).

\medskip\noindent
Conversely, suppose that (1) holds. Assume that $x, x' \in X_{S'}$
map to the same point $s' \in S'$. Choose a commutative diagram
$$
\xymatrix{
K & \kappa(x) \ar[l] \\
\kappa(x') \ar[u] & \kappa(s') \ar[l] \ar[u]
}
$$
of fields. By Schemes, Lemma \ref{schemes-lemma-characterize-points}
we get two morphisms $a, a' : \Spec(K) \to X_{S'}$. One corresponding
to the point $x$ and the embedding $\kappa(x) \subset K$ and
the other corresponding to the  point $x'$ and the embedding
$\kappa(x') \subset K$. Also we have $f' \circ a = f' \circ a'$.
Condition (1) now implies that the compositions of $a$ and $a'$ with
$X_{S'} \to X$ are equal. Since $X_{S'}$ is the fibre product
of $S'$ and $X$ over $S$ we see that $a = a'$. Hence $x = x'$.
Thus (1) implies (2).

\medskip\noindent
If there are two different points $x, x' \in X$ mapping to the same point of $s$
then (2) is violated.
If for some $s = f(x)$, $x \in X$ the field extension
$\kappa(x)/\kappa(s)$ is not purely inseparable, then
we may find a field extension $K/\kappa(s)$ such that
$\kappa(x)$ has two $\kappa(s)$-homomorphisms into $K$. By
Schemes, Lemma \ref{schemes-lemma-characterize-points} this
implies that the map
$\Mor(\Spec(K), X) \to \Mor(\Spec(K), S)$
is not injective, and hence (1) is violated.
Thus we see that the equivalent conditions (1) and (2) imply
$f$ is radicial, i.e., they imply (3).

\medskip\noindent
Assume (3). By
Schemes, Lemma \ref{schemes-lemma-characterize-points}
a morphism $\Spec(K) \to X$ is given by a pair $(x, \kappa(x) \to K)$.
Property (3) says exactly that associating to the pair
$(x, \kappa(x) \to K)$ the pair $(s, \kappa(s) \to \kappa(x) \to K)$
is injective. In other words (1) holds. At this point we know that
(1), (2) and (3) are all equivalent.

\medskip\noindent
Finally, we prove the equivalence of (4) with (1), (2) and (3).
A point of $X \times_S X$ is given by a quadruple
$(x_1, x_2, s, \mathfrak p)$, where $x_1, x_2 \in X$,
$f(x_1) = f(x_2) = s$ and
$\mathfrak p \subset \kappa(x_1) \otimes_{\kappa(s)} \kappa(x_2)$
is a prime ideal, see
Schemes, Lemma \ref{schemes-lemma-points-fibre-product}.
If $f$ is universally injective, then
by taking $S'=X$ in the definition of universally injective,
$\Delta_{X/S}$ must be surjective since it is a section of
the injective morphism
$X \times_S X  \longrightarrow X$.
Conversely, if
$\Delta_{X/S}$ is surjective, then always $x_1 = x_2 = x$ and there
is exactly one such prime ideal $\mathfrak p$, which means that
$\kappa(s) \subset \kappa(x)$ is purely inseparable.
Hence $f$ is radicial.
Alternatively, if $\Delta_{X/S}$ is surjective,
then for any $S' \to S$ the base
change $\Delta_{X_{S'}/S'}$ is surjective which implies that $f$
is universally injective. This finishes the proof of the lemma.
\end{proof}

\end{document}
